% Auto-generated by code/ — DO NOT EDIT BY HAND.
\begin{table}[htbp]
\centering
\setlength{\tabcolsep}{5pt}
\small
\caption{Population-weighted Gini of regional luminosity, 2014--2024.}
\label{tab:gini_main}
\begin{tabular}{lrrrrccr}
\toprule
Panel & Regions & \shortstack{Gini\\2014} & \shortstack{Gini\\2024} & \shortstack{Change\\in Gini} & 95\% CI & Boot. $p$ & Trend \\
\midrule
ASEAN (ten member states) & 278 & 0.584 & 0.528 & -0.057 & [-0.072, -0.042] & $<$0.001 & -0.0063 \\
Japan (benchmark)$^{\ddagger}$ & 47 & 0.337 & 0.349 & 0.011 & [-0.016, 0.025] & 0.239 & 0.0013 \\
ASEAN and Japan pooled & 325 & 0.575 & 0.517 & -0.058 & [-0.075, -0.040] & $<$0.001 & -0.0059 \\
\bottomrule
\end{tabular}
\\[2pt]
\parbox[t]{\linewidth}{\footnotesize \emph{Notes:} The Gini is the population-weighted Gini of $PW_{rt}$ across regions, computed on levels. The change is the 2024 value less the 2014 value. The interval and $p$-value come from a paired region bootstrap (2,000 draws, regions resampled with their full time path). ``Trend'' is the OLS slope of $G_t$ on the year; its conventional $p$-value assumes independent yearly deviations and is not reported, because the bootstrap interval is the appropriate inference. $^{\ddagger}$ Japan is a benchmark comparator, not an ASEAN member state.}
\end{table}
